\begin{lemma}[Maximum property under a corresponding coloring]
Let \(T\) be an extended Tashkinov tree with rung number \(n\), let
\(\sigma_n\) be a \(\varphi_n\bmod T_n\) coloring, and let \(T^*\) be a
displayed tree carrying an actual corresponding-ETT certificate for
\((\sigma_n,T_n)\).  If \(T\) has the maximum property under
\(\varphi_n\), then the replacement ETT certified on \(T^*\) has the
maximum property under \(\sigma_n\).
\end{lemma}
\begin{proof}
Write the base coloring sequence as
\(\varphi_0,\ldots,\varphi_n\) and the replacement sequence as
\(\sigma_0,\ldots,\sigma_n\).  The mod-coloring hypothesis is exactly
the paper's hypothesis that \(\sigma_n\) is
\(\varphi_n\bmod T_n\), in the sense of \noderef{ModColoring}.  Once the
displayed tree \(T^*\) and its actual corresponding certificate are fixed,
the universal extension clause in that hypothesis is not needed again.

By \noderef{CorrespondingExtendedTashkinovTree}, the replacement has the
same rung number \(n\); for every \(1\leq j\leq n\), its ordered ladder
tree \(T_j^*\) is exactly \(T_j\).  At every stored index its raw set
\(S_j^*\), exact connecting-edge set \(F_j^*\), and extension kind
\(\Theta_j^*\) equal the corresponding base data.  Moreover, each
\(\sigma_j\) is proper, leaves exactly the distinguished edge uncolored,
and, for \(1\leq j\leq n\), is
\((T_j,D_j,\varphi_j)\)-stable.  The base colorings at positive indices
are also proper and leave exactly the distinguished edge uncolored: this
is recorded in the positive SE, PE, and RE constructors of the certified
series in \noderef{TashkinovSeries} and
\noderef{TashkinovSeriesStep}.

We first record that the replacement stored set at each such stage is
really the same \(D_j\).  By \noderef{TashkinovStoredColors}, it is the
union of the raw sets \(S_h\) through \(j\), with the colors missing on
the stage tree deleted.  The corresponding certificate identifies every
raw \(S_h^*\) with \(S_h\) and identifies \(T_j^*\) with \(T_j\).
Its stability clause identifies, vertex by vertex on \(T_j\), the colors
missing under \(\sigma_j\) with those missing under \(\varphi_j\).
Taking their unions over the common vertex set and then the same finite-set
difference gives
\[
                         D_j^*=D_j.                 \tag{1}
\]

Unfold the two clauses of the maximum property in
\noderef{MaximumProperty}.  For the first clause, use ladder agreement at
the first index.  It gives \(T_1^*=T_1\) as ordered tree-sequences, hence
\(|V(T_1^*)|=|V(T_1)|\).  The base maximum property says that the latter
cardinality is the Tashkinov order of \(G\), so the replacement satisfies
the first clause as well.

For the second clause, fix \(i\) with \(1\leq i<n\).  Let \(\rho\) be a
proper coloring leaving exactly the distinguished edge uncolored which is
\((T_i^*,D_i^*,\sigma_i)\)-stable.  Let \(U\) be a prescribed seed for
the replacement data, and let \(U'\) be a closure of \(U\) under
\(\rho\).  By ladder agreement and (1), the assumed stability is precisely
\((T_i,D_i,\sigma_i)\)-stability.  The certificate also gives that
\(\sigma_i\) is \((T_i,D_i,\varphi_i)\)-stable.  All three colorings are
proper and leave exactly the same distinguished edge uncolored, so the
transitivity clause of \noderef{sc2} yields that \(\rho\) is
\((T_i,D_i,\varphi_i)\)-stable.

The same certificate identifies the extension kind and exact
connecting-edge set at index \(i\).  If the kind is PE, the replacement
seed is \(T_i^*=T_i\), which is exactly the base PE seed.  If the kind is
SE or RE, the replacement premise supplies an edge identity \(f_i\) with
the exact equality \(F_i^*=\{f_i\}\) and says that \(U\) is the exact
one-edge extension \(T_i^*+f_i\).  Agreement of \(F_i^*\) with \(F_i\)
and of \(T_i^*\) with \(T_i\) makes this the base SE/RE seed using the
same individual edge \(f_i\).  Thus \(U'\) is the very same candidate
closure to which the base maximum-property clause applies.  It gives
\[
                         |V(U')|\leq |V(T_{i+1})|.  \tag{2}
\]
Because \(i<n\), ladder agreement also includes the next stage and gives
\(T_{i+1}^*=T_{i+1}\).  Rewriting the right side of (2) proves the required
replacement bound
\(|V(U')|\leq |V(T_{i+1}^*)|\).

For clarity about indices, the certificate stores ladder agreement with a
zero-based index \(q<n\) representing the paper tree \(T_{q+1}\).  The
current stage \(T_i\) above uses \(q=i-1\), which exists because
\(1\leq i\); the next stage \(T_{i+1}\) uses \(q=i\), which exists because
\(i<n\).  Equality of the two rung numbers transports both finite indices
to the replacement series.  Finally, the quantified second clause of
\noderef{MaximumProperty} ranges only over \(1\leq i<n\).  It imposes no
maximality condition at \(i=n\), so no agreement or maximality assertion
for a final stage beyond this range is required.  The two clauses are now
proved for the replacement ETT.
\end{proof}
